% Cote 119: offprints, reports, courses, CV and letters (typed and printed part): archive cote 119, site page 251 (PDF page 252).
% typed letter. Draft compared: none.
% AI-transcribed by Claude Opus 5.5 (claude-opus-5-5, adaptive thinking, effort high), Message Batches API, batch r119_full (msgbatch_01DwZNm5Y7YsKBWyAN4sFs8w), 2026-11-13; not yet checked by eye.
\scanpage{119/251}
A. GROTHENDIECK

2, Avenue de Verrières

91 - \uline{MASSY}

A Monsieur Hubert Beuve-Méry

Directeur du Monde

5, Rue des Italiens

\uline{PARIS} IX \hfill Massy, le 4 Juin 1968

Cher Monsieur,

J'ai été invité par mes collègues du département de Mathématiques
de la Faculté des Sciences d'Orsay à participer aux travaux de la commission
chargée par ce département d'étudier les questions concernant la formation
et le recrutement des enseignants et chercheurs mathématiciens dans l'Univer-%
sité. Au cours de plusieurs discussions entre étudiants, assistants ou
maîtres assistants, et professeurs, il est apparu que certains vices fonda-%
mentaux du système actuel, dont à peu près tous les professeurs concernés
avaient une vision assez claire, étaient moins apparents pour les étudiants
et les assistants, qui sont pourtant directement concernés. Pour leur permet-%
tre de s'associer en pleine connaissance de cause aux propositions et aux
décisions qui doivent être élaborés à ce sujet dans un prochain avenir, il
m'a semblé qu'il était nécessaire d'informer sur ces problèmes un public
plus large que celui formé des quelques delégués d'étudiants et d'assistants
qui sont amenés à prendre part à ces discussions\hand{,} qui devrait comprendre pour
le moins l'ensemble des étudiants et enseignants en mathématiques, ou même
en sciences en général,- les problèmes soulevés étant sans doute les mêmes
(quoique probablement moins aigus) dans les autres sciences. C'est dans ce
but que je vous envoie ci-joint un exposé intitulé ``Le maître-enseignant et
le maître-chercheur dans l'Université d'aujourd'hui'', en vous suggérant la
publication de cet exposé, sous forme de un ou deux articles de fond, dans
un avenir rapproché. Je vous serais obligé de me faire connaître, dès qu'il
vous sera possible, s'il vous sera possible de publier cet exposé dans une

\mspage{251}
